\documentclass[10pt,a4paper]{amsart}
\usepackage{amsmath,amssymb,amsthm,graphicx}
\usepackage{verbatim}
\newenvironment{reference}{\comment}{\endcomment}
%\newenvironment{reference}{}{}
\newenvironment{slogan}{\comment}{\endcomment}
\newenvironment{history}{\comment}{\endcomment}
% For commutative diagrams we use Xy-pic
\usepackage[all]{xy}

% We use 2cell for 2-commutative diagrams.
\xyoption{2cell}
% Theorem environments.
%
\theoremstyle{plain}
\newtheorem{theorem}[subsection]{Theorem}
\newtheorem{proposition}[subsection]{Proposition}
\newtheorem{lemma}[subsection]{Lemma}

\theoremstyle{definition}
\newtheorem{definition}[subsection]{Definition}
\newtheorem{example}[subsection]{Example}
\newtheorem{exercise}[subsection]{Exercise}
\newtheorem{situation}[subsection]{Situation}

\theoremstyle{remark}
\newtheorem{remark}[subsection]{Remark}
\newtheorem{remarks}[subsection]{Remarks}

\numberwithin{equation}{subsection}

% Macros
%
\def\lim{\mathop{\mathrm{lim}}\nolimits}
\def\colim{\mathop{\mathrm{colim}}\nolimits}
\def\Spec{\mathop{\mathrm{Spec}}}
\def\Hom{\mathop{\mathrm{Hom}}\nolimits}
\def\Ext{\mathop{\mathrm{Ext}}\nolimits}
\def\SheafHom{\mathop{\mathcal{H}\!\mathit{om}}\nolimits}
\def\SheafExt{\mathop{\mathcal{E}\!\mathit{xt}}\nolimits}
\def\Sch{\mathit{Sch}}
\def\Mor{\mathop{\mathrm{Mor}}\nolimits}
\def\Ob{\mathop{\mathrm{Ob}}\nolimits}
\def\Sh{\mathop{\mathit{Sh}}\nolimits}
\def\NL{\mathop{N\!L}\nolimits}
\def\CH{\mathop{\mathrm{CH}}\nolimits}
\def\proetale{{pro\text{-}\acute{e}tale}}
\def\etale{{\acute{e}tale}}
\def\QCoh{\mathit{QCoh}}
\def\Ker{\mathop{\mathrm{Ker}}}
\def\Im{\mathop{\mathrm{Im}}}
\def\Coker{\mathop{\mathrm{Coker}}}
\def\Coim{\mathop{\mathrm{Coim}}}

% Boxtimes
%
\DeclareMathSymbol{\boxtimes}{\mathbin}{AMSa}{"02}

%
% Macros for moduli stacks/spaces
%
\def\QCohstack{\mathcal{QC}\!\mathit{oh}}
\def\Cohstack{\mathcal{C}\!\mathit{oh}}
\def\Spacesstack{\mathcal{S}\!\mathit{paces}}
\def\Quotfunctor{\mathrm{Quot}}
\def\Hilbfunctor{\mathrm{Hilb}}
\def\Curvesstack{\mathcal{C}\!\mathit{urves}}
\def\Polarizedstack{\mathcal{P}\!\mathit{olarized}}
\def\Complexesstack{\mathcal{C}\!\mathit{omplexes}}
% \Pic is the operator that assigns to X its picard group, usage \Pic(X)
% \Picardstack_{X/B} denotes the Picard stack of X over B
% \Picardfunctor_{X/B} denotes the Picard functor of X over B
\def\Pic{\mathop{\mathrm{Pic}}\nolimits}
\def\Picardstack{\mathcal{P}\!\mathit{ic}}
\def\Picardfunctor{\mathrm{Pic}}
\def\Deformationcategory{\mathcal{D}\!\mathit{ef}}
\usepackage{cleveref}
\crefname{originallemma}{original lemma}{original lemmas}
\Crefname{originallemma}{Original lemma}{Original lemmas}
\makeatletter
\newlabel{lemma-closed-in-quasi-compact}{{lemma-closed-in-quasi-compact}{source}{}{sourcepointer-1}{}}
\newlabel{lemma-closed-in-quasi-compact@cref}{{[originallemma][0][]lemma-closed-in-quasi-compact}{[1][0][]source}{}{}{}}
\makeatother
\crefname{originallemma}{original lemma}{original lemmas}
\Crefname{originallemma}{Original lemma}{Original lemmas}
\makeatletter
\newlabel{lemma-constructible-hausdorff-quasi-compact}{{lemma-constructible-hausdorff-quasi-compact}{source}{}{sourcepointer-2}{}}
\newlabel{lemma-constructible-hausdorff-quasi-compact@cref}{{[originallemma][0][]lemma-constructible-hausdorff-quasi-compact}{[1][0][]source}{}{}{}}
\makeatother
\crefname{originallemma}{original lemma}{original lemmas}
\Crefname{originallemma}{Original lemma}{Original lemmas}
\makeatletter
\newlabel{lemma-directed-inverse-limit-spectral-spaces}{{lemma-directed-inverse-limit-spectral-spaces}{source}{}{sourcepointer-3}{}}
\newlabel{lemma-directed-inverse-limit-spectral-spaces@cref}{{[originallemma][0][]lemma-directed-inverse-limit-spectral-spaces}{[1][0][]source}{}{}{}}
\makeatother
\crefname{originallemma}{original lemma}{original lemmas}
\Crefname{originallemma}{Original lemma}{Original lemmas}
\makeatletter
\newlabel{lemma-intersection-closed-in-quasi-compact}{{lemma-intersection-closed-in-quasi-compact}{source}{}{sourcepointer-4}{}}
\newlabel{lemma-intersection-closed-in-quasi-compact@cref}{{[originallemma][0][]lemma-intersection-closed-in-quasi-compact}{[1][0][]source}{}{}{}}
\makeatother
\title{Five equivalent descriptions of the spectral subsets of a spectral space}
\author{The Stacks Project Authors}
\date{stacks-project:topology.tex:local-snapshot-20261012; GNU FDL 1.2 or later}
\begin{document}
\maketitle
\noindent\textit{Editorial note.} All delivered mathematical text is the author's own, copied from the pinned Stacks Project chapter source sources/stacks-topology/pinned-original-source.tex; no statement or proof step was rewritten, reordered or removed.\par\medskip
\noindent\textit{Editorial note.} Preamble provenance: the Stacks Project keeps one single shared preamble for every chapter; this card takes it directly from the staged sources/stacks-topology/preamble.tex as three verbatim spans (lines 8-12, 14-18 and 154-217), so no derived file was written and nothing was added or redefined. The wrapper class and the base packages are supplied by the assembler.\par\medskip
\noindent\textit{Editorial note.} Bibliography provenance: neither the counted statement nor its proof contains a citation (this lane rescans the whole span before writing the card), so the card carries no bibliography and no entry was reconstructed.\par\medskip
\noindent\textit{Editorial note.} Selection provenance: this card was opened under the file-based control contract (CONTROL.md), with the mechanical pre-card gate run on its statement and proof interval before the card was written.\par\medskip
\noindent\textit{Allowed standard prerequisites (editorial).}\begin{itemize}
\item The statement that the constructible topology of a spectral space is Hausdorff and quasi-compact, and that a closed subset of a quasi-compact space is quasi-compact. Invoked inside the counted proof (source lines 5244-5247). Gives quasi-compactness of an intersection of constructible subsets.
\item The statement that a constructible subset meeting the closure of an irreducible subset contains an open of it. Invoked inside the counted proof (source lines 5258-5262). Produces quasi-compact opens containing the subset.
\end{itemize}\medskip
\ifdefined\hypertarget\hypertarget{sourcepointer-1}{}\fi
\noindent\textit{Source reference.} Original lemma lemma-closed-in-quasi-compact: sources/stacks-topology/pinned-original-source.tex, source lines 1879--1879. Navigation pointer only: the counted text refers to the Stacks item lemma-closed-in-quasi-compact, which is an existing labeled result of this same chapter and is not part of this card. The pointer carries the target's own label as its number and its definition line as its span; its content is not reproduced here.\par\medskip
\ifdefined\hypertarget\hypertarget{sourcepointer-2}{}\fi
\noindent\textit{Source reference.} Original lemma lemma-constructible-hausdorff-quasi-compact: sources/stacks-topology/pinned-original-source.tex, source lines 4522--4522. Navigation pointer only: the counted text refers to the Stacks item lemma-constructible-hausdorff-quasi-compact, which is an existing labeled result of this same chapter and is not part of this card. The pointer carries the target's own label as its number and its definition line as its span; its content is not reproduced here.\par\medskip
\ifdefined\hypertarget\hypertarget{sourcepointer-3}{}\fi
\noindent\textit{Source reference.} Original lemma lemma-directed-inverse-limit-spectral-spaces: sources/stacks-topology/pinned-original-source.tex, source lines 5119--5119. Navigation pointer only: the counted text refers to the Stacks item lemma-directed-inverse-limit-spectral-spaces, which is an existing labeled result of this same chapter and is not part of this card. The pointer carries the target's own label as its number and its definition line as its span; its content is not reproduced here.\par\medskip
\ifdefined\hypertarget\hypertarget{sourcepointer-4}{}\fi
\noindent\textit{Source reference.} Original lemma lemma-intersection-closed-in-quasi-compact: sources/stacks-topology/pinned-original-source.tex, source lines 1941--1941. Navigation pointer only: the counted text refers to the Stacks item lemma-intersection-closed-in-quasi-compact, which is an existing labeled result of this same chapter and is not part of this card. The pointer carries the target's own label as its number and its definition line as its span; its content is not reproduced here.\par\medskip
\begin{lemma}
\label{lemma-make-spectral-space}
Let $W$ be a subset of a spectral space $X$. The following are equivalent:
\begin{enumerate}
\item $W$ is an intersection of constructible sets and
closed under generalizations,
\item $W$ is quasi-compact and closed under generalizations,
\item there exists a quasi-compact subset $E \subset X$ such that
$W$ is the set of points specializing to $E$,
\item $W$ is an intersection of quasi-compact open subsets,
\item
\label{item-intersection-quasi-compact-open}
there exists a nonempty set $I$ and quasi-compact opens
$U_i \subset X$, $i \in I$
such that $W = \bigcap U_i$ and for all $i, j \in I$ there exists a
$k \in I$ with $U_k \subset U_i \cap U_j$.
\end{enumerate}
In this case we have (a) $W$ is a spectral space, (b) $W = \lim U_i$
as topological spaces, and (c) for any open $U$ containing $W$
there exists an $i$ with $U_i \subset U$.
\end{lemma}

\begin{proof}
Let $W \subset X$ satisfy (1). Then $W$ is closed in the constructible
topology, hence quasi-compact in the constructible topology (by
Lemmas \ref{lemma-constructible-hausdorff-quasi-compact} and
\ref{lemma-closed-in-quasi-compact}), hence quasi-compact in the topology
of $X$ (because opens in $X$ are open in the constructible topology). Thus
(2) holds.

\medskip\noindent
It is clear that (2) implies (3) by taking $E = W$.

\medskip\noindent
Let $X$ be a spectral space and let $E \subset W$ be as in (3).
Since every point of $W$ specializes to a point of $E$ we see that
an open of $W$ which contains $E$ is equal to $W$. Hence since $E$
is quasi-compact, so is $W$.
If $x \in X$, $x \not \in W$, then $Z = \overline{\{x\}}$ is
disjoint from $W$. Since $W$ is quasi-compact we can find a
quasi-compact open $U$ with $W \subset U$ and $U \cap Z = \emptyset$.
We conclude that (4) holds.

\medskip\noindent
If $W = \bigcap_{j \in J} U_j$ then setting $I$ equal to the set of
finite subsets of $J$ and $U_i = U_{j_1} \cap \ldots \cap U_{j_r}$
for $i = \{j_1, \ldots, j_r\}$ shows that (4) implies (5). It is immediate
that (5) implies (1).

\medskip\noindent
Let $I$ and $U_i$ be as in (5).
Since $W = \bigcap U_i$ we have $W = \lim U_i$ by the universal property
of limits. Then $W$ is a spectral space by
Lemma \ref{lemma-directed-inverse-limit-spectral-spaces}.
Let $U \subset X$ be an open neighbourhood of $W$.
Then $E_i = U_i \cap (X \setminus U)$ is a family of constructible
subsets of the spectral space $Z = X \setminus U$
with empty intersection. Using that the spectral topology on $Z$
is quasi-compact (Lemma \ref{lemma-constructible-hausdorff-quasi-compact})
we conclude from
Lemma \ref{lemma-intersection-closed-in-quasi-compact}
that $E_i = \emptyset$ for some $i$.
\end{proof}

\end{document}
